% !TeX root = ../main.tex

\chapter{与申万研报的差异归因及计算合理性证据}
\section{收益比较必须先对齐区间}
原研报[1]第6页报告截至2015年7月31日的2938个交易日、累计收益4176.87\%、年化36.74\%；本研究E的26.17\%覆盖到2026年9月。二者时期不同，不能把差额全部解释为复现错误。当前B在147个月的年化为38.86\%，实际高于36.74\%，但也不因此认定复现完全正确。

\begin{table}[htbp]
\centering
\caption{原研报与本研究收益区间的区别}
\renewcommand{\arraystretch}{1.18}
\small
\setlength{\tabcolsep}{4.5pt}
\begin{tabularx}{\textwidth}{@{}p{2.9cm}*{3}{C}@{}}
\toprule
对象与口径 & 评价月份 & 年化收益/\% & 累计收益/\% \\
\midrule
申万原文 & 原文称2003年至2015-07 & 36.74 & 4176.87 \\
既定基础复现 & 2003-04至2015-06 & 38.86 & 5481.49 \\
末组延持至7月 & 2003-04至2015-07 & 37.31 & 4891.07 \\
年初现金并延持7月 & 2003-01至2015-07 & 36.44 & 4891.07 \\
原规则完整延展 & 2003-04至2026-09 & 26.17 & 23457.11 \\
\bottomrule
\end{tabularx}
\par\smallskip
\begin{minipage}{\textwidth}\small 注：前两行不是完全同日同频比较。末组延持保留原截止前49个信号，不加入2015年6月信号，仅延长末组至7月；年初现金把1—3月未建仓收益置0。两行是起止月份敏感性，不是申万日频重建。\end{minipage}
\end{table}

终点和年化分母即可将38.86\%改变为37.31\%或36.44\%，但延持后的累计收益4891.07\%仍与4176.87\%不同。原文累计财富与年化反推约12.003年，而未披露全部年化日数约定；本文不据此认定原报告计算错误，也不把图中首个刻度当作已认证的建仓日。

\section{其他CSMAR字段与实施规则的独立对照}
本轮以现有CSMAR数据独立构造17组明确口径，在指定区间保存73行绩效。它们包括基准、单项敏感性、实施联合对照和错误字段反例，不是17个独立投资模型，也未进行组合参数寻优。表15列示主要对照，其余结果在配套分析文档保存。

\begin{table}[htbp]
\centering
\caption{统一区间后的主要口径敏感性}
\renewcommand{\arraystretch}{1.18}
\small
\setlength{\tabcolsep}{4.5pt}
\begin{tabularx}{\textwidth}{@{}p{3.35cm}*{2}{C}@{}}
\toprule
实现对照 & 147个月年化/\% & 282个月年化/\% \\
\midrule
既定月频基准 & 38.86 & 26.17 \\
行情表总市值 & 38.17 & 24.39 \\
归母利润计算PE & 44.63 & 26.74 \\
数据库同比独立历史 & 41.28 & 27.52 \\
单季同比成长 & 27.96 & 16.25 \\
五个完整年度历史 & 36.22 & 24.17 \\
财务固定滞后一季 & 18.62 & 13.39 \\
季度内权重漂移 & 38.73 & 25.82 \\
\bottomrule
\end{tabularx}
\end{table}

行情表总市值替代后原区间38.17\%，与基准较接近；归母利润计算PE为44.63\%，并未自然缩小差异。现成PE1/PE2在亏损时为空，需与同样采用正PE内排名的对照比较，不能将负PE占位变化混同于利润定义效应。

CSMAR字段字典明确F081001B是单季度环比、F081002B才是同比。原程序自行计算累计净利润同比，不存在误用字段A的问题。本轮另从04表自身建立正确同比历史，而非先丢掉03表没有的历史报告，年化为41.28\%；仍未解决末期4只股票缺失。

固定滞后一季使原区间年化降为18.62\%，表明财务期安排会改变表现，但一季滞后不等于原研报发布日算法。缺失收益按账面不变处理为38.74\%，再允许季度内权重漂移为38.73\%，本组机械收益实现影响较小。五个完整年度历史得到36.22\%，却在2015年3月选出20只、仅重合4只；把费用改到0.50\%也会接近原年化，却违背原费率。这些反例说明不能以“接近36.74\%”选择算法。

\begin{figure}[htbp]
\centering
\includegraphics[width=\linewidth]{fig14}
\caption{原区间单项口径验证}
\par\smallskip
\begin{minipage}{0.97\linewidth}\small 注：各柱均为相同147个月；较接近原报告数值不代表持仓一致。资料来源：CSMAR，作者计算。\end{minipage}
\end{figure}

\section{末期持仓差异可以定位到历史资料}
原研报第12页给出2015年5月25日调仓、7月31日持有8只。当前2015年3月信号为9只，重合4只；改用5月末市值配一季报可选出8只，仍仅重合4只。因此，数量相同不能替代名单核对。

\begin{table}[htbp]
\centering
\caption{未匹配四只股票的历史成长条件}
\renewcommand{\arraystretch}{1.18}
\small
\setlength{\tabcolsep}{4.5pt}
\begin{tabularx}{\textwidth}{@{}p{2.45cm}*{3}{C}@{}}
\toprule
股票 & 五年窗口财报数 & 正成长率数 & 缺失同比/期 \\
\midrule
康力电梯 & 20 & 17 & 3 \\
海康威视 & 20 & 16 & 4 \\
老板电器 & 18 & 14 & 4 \\
奥瑞金 & 10 & 6 & 4 \\
\bottomrule
\end{tabularx}
\end{table}

四只股票当期成长和PE名次均可通过原条件，阻碍来自历史资料：康力电梯、海康威视分别有3、4期缺少去年同期利润；老板电器、奥瑞金既缺报告期，也有4期缺同比基数。这些记录不应被说成实际负成长。04表正确同比补充后仍有缺失；05公告利润也未覆盖足够的上市前季度。

\begin{figure}[htbp]
\centering
\includegraphics[width=\linewidth]{fig15}
\caption{研报最后8只股票的条件核查}
\par\smallskip
\begin{minipage}{0.97\linewidth}\small 注：条件按当前CSMAR资料与2015年3月信号计算，不是对2015年5月25日日频价格的精确复原。资料来源：CSMAR，作者计算。\end{minipage}
\end{figure}

证据支持“当前历史覆盖会排除研报中的4只”，但不能进一步证明原作者忽略缺失、采用年频历史或拥有何种不同数据库。合理的下一步是核查CSMAR内上市前比较期财务及原文的数据说明，不是补零、豁免条件或强行加入这4只。

\section{调仓算法差异与计算证据的边界}
原研报第6页使用全市场最晚发布日与报告期后90个交易日孰早；本研究使用季度末报告期映射。这是已确认的复现近似，可能同时改变信号、成交价格和持有月份。按当前05表及所查代码范围计算，2015年一季报最晚有效公告为5月4日，原报告调仓为5月25日；共同日期仍不一致，不能直接硬填成原日期。

本轮没有2003—2015年个股日频成交重建，也未恢复当时财务修订版本。故本文证明的是既定月频定义下的计算可复核性，而非与原日频策略逐日等价。

\begin{table}[htbp]
\centering
\caption{程序与计算合理性的证据链及适用范围}
\renewcommand{\arraystretch}{1.18}
\small
\setlength{\tabcolsep}{4.5pt}
\begin{tabularx}{\textwidth}{@{}p{2.35cm}LL@{}}
\toprule
核验层次 & 实际证据 & 可以支持的结论 \\
\midrule
输入与变量 & CSMAR字段字典、单位、主键及季度累计利润核查 & 指标定义及连接具有可追溯依据 \\
独立基准重建 & 独立Python逐股重建E的560条持仓；282个月最大差约$5\times10^{-16}$ & 当前基准R实现与独立算法一致 \\
公式交叉复算 & R直接复算17组共同147个月收益与回撤，误差小于$10^{-10}$ & 实验汇总不是由单一程序自证 \\
四路径机制核验 & 初始名单、排名、三因子得分及过去收益仓位逐项检查 & 改进路径按已披露规则执行 \\
时点与资料边界 & 原调仓日及4只股票历史缺口尚未重建一致 & 不能证明申万日频等价或交易可实现性 \\
\bottomrule
\end{tabularx}
\end{table}

从方法上，季度财务构造成长与估值、月收益评价持仓，是与本项目月频近似目标相适应的设计；市值、利润、同比基数和收益单位均有定义，分母非正及缺失的处理也已公开。其合理性是“在明示目标和假设下可解释”，不是“适用于所有交易问题”。

从实现上，独立重建与R结果在持仓及浮点精度内一致，另有不同语言的公式复算和风控序列核查，足以支持“目前未发现既定规则实现中的重大算术或连接不一致”。独立程序也可能共享同一原始数据偏差，因此仍不能推出数据绝对正确、历史可交易或存在稳健超额收益。程序一致性与投资有效性必须分别评价。

\section{差异证据应按确定程度组织}
第一类是可直接确认的定义差异，包括本研究固定季度末、原研报发布日规则，以及实际收益终点分别为6月与7月。它们已由原文和源程序对应确认，不需要依靠收益接近程度来推测。第二类是当前数据支持的排除机制：末期四只股票存在历史同比或报告期缺口，按照本项目20/20条件必然无法入选。它说明本地实现为什么排除它们，但没有证明原作者采用了相同的历史覆盖或同一种缺失处理。

第三类是敏感性结果。归母利润、现成PE、财务滞后或完整年度历史等替代定义改变收益，说明输出对口径具有依赖性；它们不能反向认证哪一个定义就是原报告真正使用的版本。尤其是五个完整年度历史的年化36.22\%接近36.74\%，但末期名单仍明显不同，构成“只对收益率”的明确反例。

第四类仍待查明，包括原始数据库历史版本、精确首笔建仓日、年化日数约定及逐笔交易实现。当前证据不能给出各原因对最终差额的可加总百分比分解：某一个变化还会改变排名、名单及后续复利，多个变化并非线性叠加。本文因此使用定义对照、单项敏感性及逐股诊断共同解释，而不把CAGR差机械分摊成若干所谓贡献率。

\section{如何判断“计算合理”而不陷入自我证明}
输入层需要明确字段、单位、主键和财报期间，公式层需要说明利润同比和TTM如何构造，实施层需要检查排名、历史条件及名单，结果层需要从逐月收益复算复利和回撤。只有最终收益接近原文，无法替代前面各层检查。反过来，程序内部一致也不能认证数据供应商每条记录或公告时点都正确。

本研究独立Python从CSMAR重建基准名单和月收益，再用R展开公式复算17组实验指标。这比重复调用同一个函数多一层独立性；但两种语言共享原始资料和部分规则解释，仍可能共享数据偏差。本文据此支持的是既定月频模型具有可追溯定义且实现交叉一致，不声称已经消除原研报与本研究的一切差异。
